\documentclass[11pt]{article}
\usepackage[margin=1in]{geometry}
\usepackage{amsmath,amssymb,amsthm}
\usepackage{hyperref}

\newtheorem{theorem}{Theorem}
\newtheorem{lemma}{Lemma}
\newtheorem{corollary}{Corollary}
\newtheorem{remark}{Remark}

\newcommand{\Xclass}{\mathcal X}
\newcommand{\rhoM}{\rho_M}
\newcommand{\eps}{\varepsilon}

\title{An Orlicz-space characterization for the class \texorpdfstring{$\mathcal X$}{X}}
\author{Run fa\_banach\_001, agent lane 17}
\date{2026-07-03}

\begin{document}
\maketitle

\section*{Source and target}

The source paper is Astashkin--Lykov--Milman, \emph{Majorization revisited:
Comparison of norms in interpolation scales}, arXiv:2107.11854 \cite{ALM}.
In Section 6 the authors define the class $\Xclass$ of rearrangement invariant
spaces $X$ on $(0,\infty)$ with the following property. If a pair $(f,g)$
satisfies the Nazarov--Podkorytov single-crossing condition
\[
 f^*(t) \ge g^*(t)\quad (t\le t_0),\qquad
 f^*(t) \le g^*(t)\quad (t\ge t_0),
\]
and if $\|f\|_X\ge \|g\|_X$, then
\[
 \|f^*\chi_{(0,s)}\|_X\ge \|g^*\chi_{(0,s)}\|_X
 \qquad\text{for every }s>0.
\]
They state that a complete characterization of $\Xclass$ is open. They prove
that Lorentz spaces $\Lambda_r(\varphi)$ belong to $\Xclass$, that the
Marcinkiewicz spaces $M(t^\alpha)$ do not, and that a $\Delta_2$ Orlicz space
$L_M$ belongs to $\Xclass$ under a monotonicity assumption on
$M'(\eps x)/M'(x)$.

\section*{Result}

This packet gives a necessary and sufficient condition for the Orlicz subcase
under the same $\Delta_2$ modular regime used by the source paper.

Let $M$ be a strictly increasing Orlicz function and, for $0<\eps<1$, put
\[
 Q_\eps(u,v)=
 \frac{M(\eps v)-M(\eps u)}{M(v)-M(u)},\qquad 0<u<v.
\]

\begin{theorem}
\label{thm:orlicz-characterization}
Let $M$ be a strictly increasing $\Delta_2$ Orlicz function on $[0,\infty)$.
Then $L_M(0,\infty)$, equipped with the Luxemburg norm, belongs to
$\Xclass$ if and only if, for every $0<\eps<1$ and every $0<u<v<w$,
\[
 Q_\eps(u,v)\ge Q_\eps(v,w).
\tag{SC}
\label{eq:SC}
\]
\end{theorem}

Thus the source paper's differentiable condition is not merely a convenient
sufficient condition; after replacing derivatives by secants it is the exact
Orlicz-side dividing line.

\section*{Proof intuition}

The proof is a measure-ratio version of the source paper's Orlicz argument.
For fixed $\eps$, compare the Stieltjes measures
\[
 dM_\eps(y):=d(M(\eps y)),\qquad dM(y).
\]
Condition \eqref{eq:SC} says that the ratio $dM_\eps/dM$, interpreted through
interval averages, decreases as the level interval moves to the right.

The $NP$ crossing makes the distribution-function difference have one sign:
the $g$-excess lives below the crossing level and the $f$-excess lives above
it. If a high-scale modular comparison favors $g$, the decreasing ratio
condition forces the lower-scale modular comparison to favor $g$ as well.
This is exactly the step needed to rule out a truncation failure.

Conversely, if \eqref{eq:SC} fails, one can concentrate the lower-level
excess and higher-level excess on two adjacent intervals where the ratio is in
the wrong order. A two-step prefix plus a common low tail then gives an
explicit $NP$ pair with $\|f\|_{L_M}\ge\|g\|_{L_M}$ but with a failed prefix
comparison.

\section*{The interval-ratio lemma}

Let $\mu$ and $\mu_\eps$ be the Stieltjes measures determined on intervals by
\[
 \mu((a,b])=M(b)-M(a),\qquad
 \mu_\eps((a,b])=M(\eps b)-M(\eps a).
\]

\begin{lemma}
\label{lem:averaging}
Assume \eqref{eq:SC}. Fix $\tau>0$ and $0<\eps<1$. If $A$ and $B$ are
nonnegative measurable functions with $A$ supported in $(0,\tau)$ and $B$
supported in $(\tau,\infty)$, and if all displayed integrals are finite, then
\[
 \int A\,d\mu>\int B\,d\mu
 \quad\Longrightarrow\quad
 \int A\,d\mu_\eps>\int B\,d\mu_\eps .
\]
\end{lemma}

\begin{proof}
Condition \eqref{eq:SC} implies the following separated-interval form: whenever
intervals $I$ and $J$ satisfy $\sup I\le \inf J$ and $\mu(I),\mu(J)>0$,
\[
 \frac{\mu_\eps(I)}{\mu(I)}\ge
 \frac{\mu_\eps(J)}{\mu(J)}.
\tag{1}
\label{eq:separated-intervals}
\]
Indeed, the ratio over a union of adjacent intervals is a weighted average of
the ratios over the pieces, so the three-point condition propagates to
arbitrary intervals by subdivision.

By finite simple-function approximation, \eqref{eq:separated-intervals}
extends from intervals to measurable sets on opposite sides of $\tau$, and
then by the layer-cake formula to nonnegative functions. Thus
\[
 \frac{\int A\,d\mu_\eps}{\int A\,d\mu}
 \ge
 \frac{\int B\,d\mu_\eps}{\int B\,d\mu},
\]
with the evident convention when one denominator is zero. Since
$M$ is strictly increasing, the relevant ratios are positive. The strict
inequality $\int A\,d\mu>\int B\,d\mu$ therefore gives the claim.
\end{proof}

\section*{Sufficiency}

We prove that \eqref{eq:SC} implies $L_M\in\Xclass$. Write
\[
 \rhoM(h)=\int_0^\infty M(|h(t)|)\,dt
\]
for the Orlicz modular. Under $\Delta_2$, the standard Luxemburg modular
normalization gives $\|h\|_{L_M}=1$ iff $\rhoM(h)=1$ for nonzero $h$ in
$L_M$, and modular convergence is compatible with monotone truncations. As in
the source paper's proof of its Orlicz theorem, it is therefore enough to
argue first for continuous decreasing rearrangements with finite support and
then pass to limits.

Let $(f,g)$ satisfy the $NP$ condition and $\|f\|_{L_M}\ge \|g\|_{L_M}$. Scale
so that $\|f\|_{L_M}=1$. Then $\rhoM(f)=1$ and $\rhoM(g)\le 1$.

Suppose, toward a contradiction, that the defining truncation inequality fails
at some $s_0$. It cannot fail for $s_0\le t_0$, so $s_0>t_0$. Let
$F=f^*\chi_{(0,s_0)}$ and $G=g^*\chi_{(0,s_0)}$. Since
$\|F\|_{L_M}<\|G\|_{L_M}\le 1$, choose $\lambda$ with
\[
 \|F\|_{L_M}<\lambda<\|G\|_{L_M}\le 1.
\]
Set $a=1/\lambda>1$ and $\eps=\lambda$. Then
\[
 \int_0^{s_0} M(aG(t))\,dt>1\ge
 \int_0^{s_0} M(aF(t))\,dt.
\tag{2}
\label{eq:high-scale-prefix}
\]

Let $\Phi=aF$ and $\Psi=aG$. The $NP$ crossing implies that, for some level
$\tau$, the distribution-function difference has the sign pattern
\[
 \lambda_\Psi(y)-\lambda_\Phi(y)\ge0\quad (0<y<\tau),\qquad
 \lambda_\Phi(y)-\lambda_\Psi(y)\ge0\quad (y>\tau).
\]
By the layer-cake/Stieltjes representation of the modular, inequality
\eqref{eq:high-scale-prefix} says
\[
 \int_0^\tau (\lambda_\Psi-\lambda_\Phi)\,d\mu
 >
 \int_\tau^\infty (\lambda_\Phi-\lambda_\Psi)\,d\mu .
\]
Applying Lemma \ref{lem:averaging} to
$A=(\lambda_\Psi-\lambda_\Phi)_+$ below $\tau$ and
$B=(\lambda_\Phi-\lambda_\Psi)_+$ above $\tau$ yields
\[
 \int_0^{s_0} M(G(t))\,dt>
 \int_0^{s_0} M(F(t))\,dt .
\tag{3}
\label{eq:low-scale-prefix}
\]
Since $s_0>t_0$, the $NP$ condition gives $g^*(t)\ge f^*(t)$ for
$t>s_0$. Hence the tail modular of $g$ is at least the tail modular of $f$.
Combining this tail comparison with \eqref{eq:low-scale-prefix} gives
$\rhoM(g)>\rhoM(f)$, contradicting $\rhoM(g)\le\rhoM(f)=1$. Therefore no
truncation failure is possible, and $L_M\in\Xclass$.

\section*{Necessity}

It remains to show that failure of \eqref{eq:SC} forces
$L_M\notin\Xclass$. Suppose that for some $0<\eps<1$ and $0<d<a<b$,
\[
\frac{M(\eps a)-M(\eps d)}{M(a)-M(d)}
<
\frac{M(\eps b)-M(\eps a)}{M(b)-M(a)}.
\tag{4}
\label{eq:bad-secant}
\]
Equivalently, there is $\theta>0$ such that
\[
 \frac{M(b)-M(a)}{M(a)-M(d)}
 < \theta <
 \frac{M(\eps b)-M(\eps a)}
      {M(\eps a)-M(\eps d)} .
\tag{5}
\label{eq:theta}
\]
The left inequality gives
\[
 M(b)+\theta M(d)<(1+\theta)M(a).
\tag{6}
\label{eq:large-scale}
\]
Choose $B>0$ so that
\[
 B\bigl(M(b)+\theta M(d)\bigr)<1
 <B(1+\theta)M(a),
\tag{7}
\label{eq:B-choice}
\]
and put $C=\theta B$ and $s_0=B+C$.

On $(0,s_0)$ define decreasing step functions
\[
 U_0(t)=
 \begin{cases}
 b,&0<t<B,\\
 d,&B<t<s_0,
 \end{cases}
 \qquad
 V_0(t)=a\quad(0<t<s_0).
\]
Then \eqref{eq:B-choice} says
\[
 \int_0^{s_0}M(U_0(t))\,dt<1<
 \int_0^{s_0}M(V_0(t))\,dt.
\tag{8}
\label{eq:prefix-large}
\]
The right inequality in \eqref{eq:theta} gives
\[
 B M(\eps b)+C M(\eps d)
 >(B+C)M(\eps a).
\tag{9}
\label{eq:small-scale}
\]
Because $\eps<1$ and $M$ is increasing, the left side of
\eqref{eq:small-scale} is strictly smaller than
$B M(b)+C M(d)$, hence is $<1$ by \eqref{eq:B-choice}. Therefore
\[
 T=\frac{1-\int_0^{s_0}M(\eps U_0(t))\,dt}{M(\eps d)}
\]
is positive. Extend by a common low tail:
\[
 U(t)=
 \begin{cases}
 U_0(t),&0<t<s_0,\\
 d,&s_0<t<s_0+T,\\
 0,&t>s_0+T,
 \end{cases}
 \qquad
 V(t)=
 \begin{cases}
 V_0(t),&0<t<s_0,\\
 d,&s_0<t<s_0+T,\\
 0,&t>s_0+T.
 \end{cases}
\]
Finally set $f=\eps U$ and $g=\eps V$.

These functions are already decreasing rearrangements. With $t_0=B$ they
satisfy the source paper's $NP$ condition: $f^*\ge g^*$ on $(0,B)$ and
$f^*\le g^*$ on $(B,\infty)$.

At scale $1$, the choice of $T$ gives
\[
 \int_0^\infty M(f(t))\,dt=1,
\]
while \eqref{eq:small-scale} gives
\[
 \int_0^\infty M(g(t))\,dt<1.
\]
Thus $\|f\|_{L_M}=1>\|g\|_{L_M}$. On the prefix $(0,s_0)$,
\eqref{eq:prefix-large} says
\[
 \int_0^{s_0}M(f(t)/\eps)\,dt<1<
 \int_0^{s_0}M(g(t)/\eps)\,dt,
\]
so
\[
 \|f^*\chi_{(0,s_0)}\|_{L_M}<\eps<
 \|g^*\chi_{(0,s_0)}\|_{L_M}.
\]
This violates the defining implication for $\Xclass$. Hence
$L_M\notin\Xclass$.

The sufficiency and necessity parts prove Theorem
\ref{thm:orlicz-characterization}.

\section*{Consequences}

\begin{corollary}
The source paper's differentiable condition implies \eqref{eq:SC}. In
particular, Theorem \ref{thm:orlicz-characterization} recovers its positive
Orlicz theorem.
\end{corollary}

\begin{proof}
If $M$ is differentiable and $x\mapsto M'(\eps x)/M'(x)$ is nonincreasing,
then
\[
 Q_\eps(u,v)
 =\frac{\int_u^v \eps M'(\eps x)\,dx}{\int_u^v M'(x)\,dx}
 =\frac{\int_u^v \eps\,M'(\eps x)/M'(x)\,dM(x)}
        {\int_u^v dM(x)}.
\]
Thus $Q_\eps(u,v)$ is the $dM$-average over $(u,v)$ of a nonincreasing
function. Such averages decrease as the interval moves to the right, giving
\eqref{eq:SC}.
\end{proof}

\begin{corollary}
\label{cor:decreasing-power}
Let $0<p<q$ and define
\[
 N_{p,q}(u)=
 \begin{cases}
   u^q, & 0\le u\le 1,\\
   u^p, & u\ge 1,
 \end{cases}
 \qquad
 M_{p,q}(u)=\int_0^u N_{p,q}(v)\,dv .
\]
Then $M_{p,q}$ is a $\Delta_2$ Orlicz function and
$L_{M_{p,q}}(0,\infty)\notin\Xclass$.
\end{corollary}

\begin{proof}
The function $N_{p,q}$ is increasing and continuous, so $M_{p,q}$ is an
Orlicz function, and its two polynomial growth regimes imply $\Delta_2$.
Fix $0<\eps<1$ and choose $0<d<a<1$. Then
\[
\frac{M_{p,q}(\eps a)-M_{p,q}(\eps d)}
     {M_{p,q}(a)-M_{p,q}(d)}
=\eps^{q+1}.
\]
As $b\to\infty$,
\[
\frac{M_{p,q}(\eps b)-M_{p,q}(\eps a)}
     {M_{p,q}(b)-M_{p,q}(a)}
\longrightarrow \eps^{p+1}.
\]
Since $0<p<q$ and $0<\eps<1$, $\eps^{q+1}<\eps^{p+1}$. Hence
\eqref{eq:SC} fails for sufficiently large $b$, and the characterization
gives $L_{M_{p,q}}\notin\Xclass$.
\end{proof}

\section*{Concrete sanity check}

The script \texttt{code/verify\_concrete.py} checks the explicit decreasing
power instance $p=1$, $q=2$, $\eps=1/2$, $d=1/4$, $a=1/2$, $b=20$. It prints
\[
0.125=
\frac{M(\eps a)-M(\eps d)}{M(a)-M(d)}
<
\frac{M(\eps b)-M(\eps a)}{M(b)-M(a)}
\approx 0.2494,
\]
then constructs admissible $B,C,T$ and verifies the prefix and global modular
inequalities used in the necessity construction.

\section*{Limitations and novelty check}

This packet does not solve the source paper's complete problem for all
rearrangement invariant spaces. It gives a full characterization of the
$\Delta_2$ Orlicz subcase and turns the source paper's differentiable
sufficient condition into an exact secant-ratio criterion.

On 2026-07-03, web searches for exact and close phrases around
\emph{Majorization revisited}, $\Xclass$, Orlicz spaces, Nazarov--Podkorytov,
and the ratio $M'(\eps x)/M'(x)$ found the source arXiv page but no later
answer or matching Orlicz characterization. Novelty confidence is moderate,
pending expert literature review.

\section*{Human-review recommendation}

The main points to check are Lemma \ref{lem:averaging}, the use of the
standard $\Delta_2$ Luxemburg modular normalization, and the reduction from
continuous decreasing rearrangements to the general Orlicz-space case. The
necessity direction is explicit and finite-dimensional.

\begin{thebibliography}{9}
\bibitem{ALM}
S. V. Astashkin, K. V. Lykov, and M. Milman,
\emph{Majorization revisited: Comparison of norms in interpolation scales},
arXiv:2107.11854, 2021.
\end{thebibliography}

\end{document}
